\section{Small extension fields and higher-genus pipelines}\label{sec:extension}

This work is where the programme most clearly explores known
non-generic asymptotic routes rather than trying to manufacture prime-field
structure. The results are interesting precisely because complete costs
can be measured against the same problem. They remain separate from
deployed prime-field curves and prime-degree binary challenges.

\subsection{The residual-walk and parity ledgers}

The $\F_{p^3}$ harness studies subspace factor bases, residual walks,
large primes, GLV symmetries, decomposition oracles and sparse matrices.
Its corrections show that an oracle that reduces the number of walked
residuals can make the method more expensive if the per-residual work
increases. One cited example reduces a headline walked count by around
two orders of magnitude while total cost worsens by about 4.6 times.
\src{RESIDUAL}

The same record prices linear algebra that earlier rounds omitted. The
linear-algebra growth then dominates a relation-phase exponent and
changes the claimed crossover. This provides an unusually concrete
negative result: the bottleneck is named and priced, rather than the
route being declared unsuccessful because it did not immediately beat
rho. The parity ledger extends this discipline to several degrees and
decomposition choices.\src{PARITY}

\subsection{Degree four: symmetry improves a real algebraic constant}

On suitable curves over $\F_{p^4}$, a rational two-torsion point and an
appropriate coordinate make the summation system invariant under a
$D_4$ action. The unweighted invariant formulation reduces a generic
4,096-dimensional quotient to 512. Its recorded decomposition constant
at $p=269$ is about $2.54\cdot10^{10}$ base-field multiplications,
approximately 48 times below the plain formulation.\src{PRIOR}

The weighted four-variable follow-up is stronger evidence of iteration
than quoting that original factor forever. It replaces the five-variable
presentation by invariants of weights $(2,2,2,1)$ and records a
3,979-row, 4,200-column matrix rather than the previous 13,412 by
11,628 matrix. At $p=269,1033,4093$, median counted costs are around
$1.4\cdot10^9$ multiplications including lift-back and reruns at the
next degree when required. This is roughly another 18 times below the
previous invariant formulation, and about 844--863 times below the
named plain reference at those primes.\src{D4}

The study keeps exceptions. At $p=17$, one of six planted targets is
degenerate; the other five agree with brute force and MITM. The nine
generic positive-control systems have the expected quotient dimension.
The gain is a decomposition-stage constant. The study explicitly does
not measure end-to-end $S$ for that weighted formulation.

Its crossover calculations are therefore model outputs. Earlier models
moved a degree-four handover from around $2^{150}$ to $2^{135}$; the
weighted follow-up reports a further handover near $2^{118.9}$ in its
particular model. Those are not three observed ECDLP crossings. They
depend on the relation rate, cofactor, matrix model, unit conversion and
the constant measured at much smaller sizes.\src{D4}\src{LEADS18}

The subsequent complete-method experiments also show that combining
individually attractive ideas need not help. A D4-plus-double-large-prime
variant makes linear algebra cheaper but total cost worse on its
held-out panel. Stopping at the first actual dependency then repairs
part of that loss: the later complete candidate/incumbent ratio is
0.979 [0.973,0.986] over 27 held-out curves. This is progress against the
incumbent, not a claim that the weighted stage and the complete variant
were one already-measured stack.\src{ROUND1}\src{ROUND2}

\subsection{Degree five: a large stage factor with a distant model crossing}

The quintic Joux--Vitse/torsion programme implements a four-point
decomposition test over $\F_{p^5}$. Its plain and Edwards-symmetry systems
have very different costs: the cited constants are roughly
$3.0\cdot10^{10}$ and $7.3\cdot10^7$ base-field multiplications, a
recorded stage factor around 416. But a corrected exact yield census
changes the residual rate from $1/(24p)$ to $1/(192p)$, worsening the
extrapolated crossover. A four-torsion follow-up then retracts an
anticipated degree-halving lever: the available translation does not
implement the needed involution over the stated base field, and the
measured saturation costs about as much as it gains in rate.
\src{QUINTIC}

These examples explain why a several-hundred-fold algebraic improvement
can coexist with an impractical complete-method prediction. The
programme has learned both a useful symmetry implementation and the
conditions under which its hoped-for scaling argument fails. Neither
observation should be compressed into simply ``torsion works'' or
``torsion does not work''.

\subsection{Covers and genus-three index calculus}

The cover programme implements extension fields, Cantor-style divisor
arithmetic, a GHS/Joux--Vitse-type cover, a point transfer and a
decomposition test checked against a separate enumeration oracle.
On the favourable weak class, reducing the oracle constant improves the
complete IC/rho ratio as size grows. It reproduces a known attack route;
the difficult additional question is whether a generic source can find
and reach such a class at a total cost below its own rho reference.
\src{COVER}

The more recent cubic-GHS/genus-three work adds double large primes and
new matrix strategies. The first improvement round reports a complete
cost exponent about 1.385 rather than 1.910 on a common
$q=2^{10}$--$2^{18}$ range. Earlier headlines compared different ranges
and are corrected. The candidate's local exponent rises at larger
sizes, so a global finite-range fit is not evidence that it has reached
the asymptotic $4/3$ prediction.\src{ROUND1}

A separately rebuilt genus-three baseline reports IC/rho around 0.252
at $q=2^{14}$ after the study's rho normalisation. That comparator is
rho on the same Jacobian subgroup, not rho on the original elliptic
curve before transport. The timeline explicitly labels it as a
Jacobian result. An ECDLP crossover requires the source transfer and
source reference to be in the comparison.\src{TIMELINE}

The later structured-elimination and block-Wiedemann rounds reduce
complete costs further, but also identify the relation phase and matrix
growth that remain. This is the strongest recorded fitted-exponent
improvement over an incumbent among the selected complete toy pipelines.
It is appropriately described as reproducing and engineering a known
higher-genus route at toy scale, not discovering a new exponent for
general prime-field ECDLP.\src{ROUND2}
